% vo_style.tex -- shared look for the paper's data figures (Sec. 5).
% Each figure loads this automatically. Edit here to restyle every figure at once.
%
% Compile with LuaLaTeX (Overleaf: Menu -> Compiler -> LuaLaTeX).
% The paper preamble needs, once:
%   \usepackage[T1]{fontenc}   % without it LuaLaTeX sets the whole paper in Latin Modern, not Times
%   \usepackage{pgfplots}
%   \usepgfplotslibrary{groupplots}
%   \pgfplotsset{compat=1.18}
\def\vostyleloaded{}
\usetikzlibrary{calc}   % figure layouts offset axes with ($(a)+(x,y)$)

% ---- Colors: the Fig. 1 palette -------------------------------------------------------
% \providecolor keeps Fig. 1's own definitions if they are loaded first.
\providecolor{colorA}{RGB}{30,120,200}     % ego E
\providecolor{colorB}{RGB}{220,160,0}      % obstacle O
\providecolor{colorRisk}{RGB}{200,30,30}   % risk: collision / condition holds
\providecolor{colorSafe}{RGB}{44,160,44}   % safe: no collision / condition fails

% ---- Sample marks ----------------------------------------------------------------------
% Fill-only marks: the built-in mark=* and square* also stroke an outline, which is wider
% than the mark at this size.
\pgfdeclareplotmark{vosquare}{%
  \pgfpathrectangle{\pgfpoint{-\pgfplotmarksize}{-\pgfplotmarksize}}%
                   {\pgfpoint{2\pgfplotmarksize}{2\pgfplotmarksize}}\pgfusepathqfill}
\pgfdeclareplotmark{vodot}{\pgfpathcircle{\pgfpointorigin}{\pgfplotmarksize}\pgfusepathqfill}
% Circles by default. They need LuaLaTeX: pdfLaTeX's fixed memory runs out above ~3k circle
% marks per panel. Under pdfLaTeX use squares instead:  \renewcommand{\vomark}{vosquare}
\providecommand{\vomark}{vodot}
\providecommand{\vomarksize}{0.5pt}
\providecommand{\vomarkopacity}{0.6}

\pgfplotsset{
  vo/axis/.style={
    set layers, table/col sep=comma,
    tick label style={font=\scriptsize}, label style={font=\small}, title style={font=\small},
    grid=major, grid style={draw=black!10, line width=0.3pt},
  },
  vo/marginal/.style={vo/axis, xtick=\empty, ytick=\empty, grid=none},
  vo/marks/.style={only marks, mark=\vomark, mark size=\vomarksize, opacity=\vomarkopacity},
  % Sample classes. The numbered variants shade by magnitude: 0 = smallest (base color)
  % to 3 = largest (darkest).
  vo/hold/.style={vo/marks, color=colorRisk},
  vo/hold 0/.style={vo/marks, color=colorRisk!94!black},
  vo/hold 1/.style={vo/marks, color=colorRisk!81!black},
  vo/hold 2/.style={vo/marks, color=colorRisk!69!black},
  vo/hold 3/.style={vo/marks, color=colorRisk!56!black},
  vo/fail/.style={vo/marks, color=colorSafe},
  vo/fail 0/.style={vo/marks, color=colorSafe!94!black},
  vo/fail 1/.style={vo/marks, color=colorSafe!81!black},
  vo/fail 2/.style={vo/marks, color=colorSafe!69!black},
  vo/fail 3/.style={vo/marks, color=colorSafe!56!black},
  % Markers
  vo/apex/.style={only marks, mark=square*, mark size=1.8pt, color=colorB},
  vo/ego/.style={only marks, mark=*, mark size=1.8pt, color=colorA},
  vo/cpa/.style={only marks, mark=*, mark size=1.4pt, color=black},
  vo/twin hold/.style={only marks, mark=*, mark size=2.4pt, fill=colorRisk, draw=black, line width=0.5pt},
  vo/twin fail/.style={only marks, mark=*, mark size=2.4pt, fill=colorSafe, draw=black, line width=0.5pt},
  % Marginal histograms
  vo/hist/.style={draw=none, fill=black!45, fill opacity=0.65},
}

% ---- Lines and annotations (TikZ styles, for \draw and \node inside axes) ----------------
\tikzset{
  vo boundary/.style={draw=black!40, line width=0.9pt},            % a criterion's boundary
  vo tangent/.style={draw=black, line width=0.6pt, dash pattern=on 3pt off 2pt},
  vo cone/.style={draw=black, line width=0.9pt},                   % deterministic VO
  vo band/.style={fill=colorRisk, fill opacity=0.08},              % collision band
  vo rline/.style={draw=colorRisk, line width=0.6pt, dash pattern=on 3pt off 2pt},
  vo zero/.style={draw=black!45, line width=0.4pt},
  vo detline/.style={draw=black, line width=0.6pt, dash pattern=on 1pt off 1.5pt},
  vo label/.style={font=\scriptsize, fill=white, inner sep=1pt},
}
